%%%%%%%%%%%%%%%%%%%%%%%%%%%%%%%%%%%%%%%%%%%%%%%%%%%%%%%%%%%%%%%%%%%%%%%%%%%%%%%%%%%%
% Template for STAT 548 Qualifying Paper Report
% Author: Ben Bloem-Reddy <benbr@stat.ubc.ca>
% Revised: Daniel J. McDonald <daniel@stat.ubc.ca>
% Date: 26 August 2024
%%%%%%%%%%%%%%%%%%%%%%%%%%%%%%%%%%%%%%%%%%%%%%%%%%%%%%%%%%%%%%%%%%%%%%%%%%%%%%%%%%%%

% Note: You will get an empty bibliography warning when compiling until you include a citation.

\documentclass[12pt]{article}
\usepackage[commenters={JJZ}]{shortex} % replace OA with your initials
\usepackage[round]{natbib}
\usepackage[margin=2.5cm]{geometry}
\newcommand{\email}[1]{\href{mailto:#1}{#1}}

% minor adjustments to ShorTeX
\let\argmin\relax\DeclareMathOperator*{\argmin}{argmin}
\let\argmax\relax\DeclareMathOperator*{\argmax}{argmax}
\DeclareMathOperator*{\minimize}{minimize}
\DeclareMathOperator*{\maximize}{maximize}
\DeclareMathOperator{\subjto}{\ \text{subject to}\ }
\renewcommand{\top}{\mathsf{T}}
\renewcommand{\d}{\mathsf{d}}

\graphicspath{{fig/}}



%%%%%%%%%%%%%%%%%%%%%%%%%%%%%%%%%%%%%%%%%%%%%%%%%%%%%%%%%%%%%%%%%%%%%%%%%%%%%%%%%

% your title/author/date information go here
\title{\LaTeX\ Template for STAT 548 Qualifying Paper Report} % replace with your title, a meaningful title
\author{Justin J Zhang} % replace with your name
\date{\today} % replace with your submission date


% start of document
\begin{document}
\maketitle

\newpage

\section{Introduction}
Since the industrial revolution, the earth has seen a gradual increase in carbon emissions, leading to climate change and an increase in yearly natural disasters.
Calamities like flooding and forest fires are becoming increasingly common, creating a need to carefully monitor and predict these events for human safety.
An example is the increasing possibility of drought-like conditions in the UK, which is monitored through temperature and dryness metrics, which we denote $X_t, Y_T$ respectively. 
Since we only care about how these variables behave at high temperature and dryness, we can use Extreme Value Theory (EVT). 
Our quantity of interest is the \textit{return curve}, given in \cref{eq: rc}
\begin{equation} \label{eq: rc}
  RC_t(p) = \{ (x, y) \in \mathbb{R}^2 \ |\  \mathbb{P}(X_t > x, Y_t > y) = p \}
\end{equation}
This curve represents the temperature, dryness pairs that have a $p$-probability of occurence, illustrating the joint tail behaviour. 
Of particular importance is capturing the dependence relation between temperature and dryness since common sense dictates their are positively correlated.
When dependence is stationary, $RC_t(p)$ is independent of the time $t$, however, in practice data exhibits non-stationarity, meaning the return curve changes over time.
There are a few pre-existing methods that account for non-stationarity, however these all assume asymptotic dependence in the data, when there is often asymptotic independence.

\textit{Modelling non-stationarity in asymptotically independent extremes} by Murphy-Barltrop and Wadsworth, hencefoth denoted MBW, propose 
three methods to quantify dependence for arbitrary asymptotic behvaiour: quantile regression, smooth polynomials, and GAM. In the following sections, we will provide an overview of existing
relevant extreme value theory, introduce and explain the three methods of dependence estimation,touch on limitations, and provide a simulation study.

\section{Background Works}
Though there has not been any methods developed for our specific use case, there exists theory on modelling non-stationarity for univariate extremes (Youngman, 2019), bivariate extremes (REFERNCE), which are prudent to give background for.

  \subsection{Non-Stationarity in Univariate Extremes}
  In this section, let $\{ Y_t \}$ be a non-stationary process with covariates $z_t \in \mathbb{R}^g$.
  We define \textit{exceedences} to be $Y_t - u$ for some threshold $u \in [0, 1]$. In stationary settings, 
  (normalized) exceedences are modeled by the \textit{generalized pareto distribution} (GPD) through the Pickands-Balkema-de Haan theorem (Pickands, 1975). A GPD is parameterized by a scale $\tau \in (0, \infty)$ and shape $\xi \in \mathbb{R}$. Analogously, for our non-stationary process, the left side of \cref{eq: gpd} shows that exceedences are likewise approximated by a GPD, however with scale and shape parameters defined in terms of $z_t$ (Youngman, 2019). 
  These are estimated through an invertible link function $h(x)$ and general additive model (GAM).
  \begin{equation} \label{eq: gpd}
    Y_t - u \ |\ Y_t > u, z_t \sim \GP(\tau(z_t), \xi(z_t)) \qquad h(\zeta(z_t)) = \beta_0 + \sum_{k = 1}^g \sum_{p = 1}^{P_k} \beta_{kp}b_{kp}(z_{kt})
  \end{equation}
  For the GAM formulation in the right side of \cref{eq: gpd}, $b_{kp}$ are specified basis functions and $\beta_0, \beta_{kp}$ are coefficients estimated through penalized likelihood for each basis function $k = 1,\dots, g$ and basis dimension $p = 1,\dots, P_k$. 
  For example, we can take $h(x) = \log(x)$ and we would be fitting a GAM to $\exp(h(\zeta(z_t)))$ to obtain the parameter $\zeta(z_t)$.
  Note since $\tau(z_t) \in (0, \infty)$ we must have its link function $h_\tau^{-1}(x)$ strictly positive. 

  The functional representation of scale and shape parameters allows us to account for the change in distributional shape and variation of exceedences over time. 
  These changesn are caused by non-stationarity in the underluying process.
  Using GAMs ensures smoothness of the parameters over time, meaning that expected exceedances will not drastically change between time-steps.

  \subsection{Bivariate EVT}
  In this section, let $(X, Y)$ be a random vector representing a stationary distribution of a stochastic process. 
  We note that EVT can be done in higher dimensions but we will restrict our analysis in this paper to the bivariate case.
  We will assume the data exhibits standard exponential margins for $X, Y$, i.e. $f_X(z) = f_Y(z) = e^{-z}$. 
  Otherwise, we can apply the probability integral transform (APPENDIX) to achieve this. Then for extreme thresholds $u$, 
  the joint tail distribution can be modeled in terms of a \textit{slowly varying function} $L$ and \textit{coefficient of tail dependence} $\eta \in (0, 1]$ as shown in \cref{eq: stat joint} (Ledford, Tawn, 1996).
   Though this equality holds as $u \uparrow 1$, the equation is a good approximation for joint tail probability when $u$ is reasonably high, say 0.95.
  \begin{equation} \label{eq: stat joint}
    \lim_{u \to \infty}\mathbb{P}(\min(X, Y) > u) = \lim_{u \to \infty} L(e^u)e^{-\frac{u}{\eta}} \qquad \lim_{x \to \infty} \frac{L(cx)}{L(x)} = 1, \quad c > 0
  \end{equation}
  Asymptotic dependence occurs when $\lim_{u \uparrow 1} L(e^u) = 0$ and $\eta = 1$, which together imply the joint distribution at high thresholds is non-zero. 
  To illustrate \cref{eq: stat joint}, we can take the trivial example where $X, Y$ are independent, and with the standard exponential margins, we get for any $u$ that $\mathbb{P}(\min(X, Y) > u) = e^{-2u}$. 
  In this case, $L(x) = 1$ is constant and $\eta = \frac{1}{2}$.
  
  Though this model captures both asymptotic independence and dependence, for the latter it fails to capture regions where one variable is more extreme than the other. 
  To account for varying levels of extremeness, we build upon \cref{eq: stat joint} using different rays $w \in [0, 1]$ from the origin, shown in \cref{eq: stat joint lambda}. 
  The coefficient of tail dependence is replaced with the \textit{angular dependence function} $\lambda(w)$ (Wadsworth, Tawn, 2013). 
  \begin{equation} \label{eq: stat joint lambda}
    \lim_{u \to \infty}\mathbb{P}\left(\min\left(\frac{X}{w}, \frac{Y}{1 - w} \right) > u \right) = \lim_{u \to \infty} L(e^u; w)e^{-\lambda(w) u}, \qquad \lambda(w) \geq \max(w, 1 - w)
  \end{equation}

  \subsection{Angular Dependence Function}
  In this section, let $(X, Y)$ be as before. We will delve further into the theory of $\lambda(w)$, which was first introduced in (Wadsworth, Tawn, 2013) and is given in \cref{eq: stat joint lambda}. 
  The more general definition considers $\kappa(\alpha, \beta)$ as the angular dependence, however, we use the simple (normalized) formualation with $\alpha = w, \beta = 1 - w$ as the ray.
  To start, we wish to define some properties for $\lambda(w)$. We call random variable $X, Y$ \textit{positive quandrant dependent} if $\forall (x, y) \in \mathbb{R}^2, \; \mathbb{P}(X \leq x, Y \leq y) \geq \mathbb{P}(X \leq x)\mathbb{P}(Y \leq y)$
  \begin{proposition} \label{prop: lambda}
    Let $\lambda(w)$ be as defined in \cref{eq: stat joint lambda}. $w, w_1, w_2 \in [0, 1]$.
    \begin{enumerate}
      \item $\lambda(1) = \lambda(0) = 1$.
      \item $\max(w, 1 - w) \leq \lambda(w, 1 - w)$.
      \item Under positive quadrant dependence, $\lambda(w) \leq 1$.
      \item $\frac{w_1}{\lambda(w_1)} \leq \frac{w_2}{\lambda(w_2)}$ and $\frac{1 - w_1}{\lambda(w_1)} \geq \frac{1 - w_2}{\lambda(w_2)}$
    \end{enumerate}
  \end{proposition}

  Proofs provided in (APPENDIX). The first three properties give boundary conditions (corresponding to independence) and provide bounds for angular dependence.
  The fourth is used to prove theoretically viability of the ADF in MBW.
  Dependence is highest when $\lambda(w) = \max(w, 1 - w)$, and independence occurs when $\lambda(w) = 1$. In general, $\lambda(w)$ is convex with minimum at $\max(w, 1 - w)$.
  The upper bound on dependence only holds with positive quandrant dependence, which is necessary to ensure the joint tail probability is non-zero.

  In relation to EVT, considering the conditional extremal probability of (X, Y) with angular dependence, there is a simple representation analogous to the memory property of exponential, as shown in \cref{eq: cond lambda deriv}.
  \begin{equation} \label{eq: cond lambda deriv}
    \lim_{u \to \infty}\mathbb{P}\left(\min\left(\frac{X}{w}, \frac{Y}{1 - w} \right) > tu \ \bigg |\ \min\left(\frac{X}{w}, \frac{Y}{1 - w} \right) > u \right) = e^{-\lambda(w) t}, \quad t \geq 1
  \end{equation}
 
  To estimate $\lambda(w)$, MBW use the \textit{Hill Estimator}. For a given threshold $u$, $v_{(1)},\dots,v_{(k)}$ are the ordered exceedences of $u$ with $v_i = \min\left(\frac{x_i}{w}, \frac{y_i}{1 - w} \right)$. We can then define the estimate $\hat{\lambda}$ as in \cref{eq: hill}. 
  For consistency of the hill estimator, we need the number of exceedences to approach $\infty$ as the total observations tend to $\infty$, but the proportion of exceedences to tend to 0.
  In practice, this assumption can be reasonably satisfied (Resnick).
  \begin{equation} \label{eq: hill}
    \hat{\lambda} = \left( \frac{1}{k - 1}\sum_{i = 1}^{k - 1} \log \frac{v_{(i)}}{v_{(k)}} \right)^{-1}
  \end{equation}
  

\section{Main Contributions}
  We have seen models that attempt to model joint extremes under a variety of conditions. 
  In this section, we focus on modelling the specific use case of non-stationarity with asymptotic independence, which MBW argue is a novel problem, through the ADF.

  For the remainder of this paper, let $\{ X_t, Y_t\}$  be a non-stationary process with standard, exponential margins (again, these can be obtained with probability integral transform).  
  Moreover, assume that dependent on some covariates $Z_t \in \mathbb{R}^g$, that $\{ X_t, Y_t \ |\ Z_t\}$ is a stationary process.
  Let $K_{w, t} = \min\left(\frac{X_t}{w}, \frac{Y_t}{1 - w} \right)$ be the \text{mini-projection} onto ray $w$. We are looking to estimate $\lambda$ in \cref{eq: cond lambda main}.
  \begin{equation} \label{eq: cond lambda main}
    \mathbb{P}(K_{w, t} > v \ |\ K_{w, t} > u, Z_t = z_t) \approx e^{-(v - u)\lambda(w \ |\ z_t)} ,\quad v > u
  \end{equation}

  Using the hill estimator, we would get a pointwise estimate for each covariate $Z_t = z_t$ (which we condition on). 
  We often do not have enough realizations for each $Z_t = z_t$ and so the pointwise estimates will be poor. Moreover, $\lambda(w)$ will not be a smooth function.
  \cref{eq: cond lambda main} also shows that the exceedences $K_{w, t} - u_{w, t} \ |\ K_{w, t} > u_{w, t}, z_t$ are exponentially distributed with rate $\lambda(w \ |\ z_t)$ through uniqueness of CDF. 
  This representation can lead to simple likelihood estimates of the ADF but due to non-stationarity, we have an estimate for each time-step, and so the resulting parameter estimates are not consistent (APPENDIX MAYBE). 
  As such, MBW propose estimating $\lambda(w)$ with a quantile-based method, Bernstein-Bezier polynomials, and GAMs, which we introduce in this section.

  \subsection{Quantile-based Estimator}
    Our first method for estimating $\lambda(w \ |\ z_t)$ is a quantile method using \cref{eq: cond lambda main}.
    Consider extreme quantiles $q_1 < q_2 < 1$ (generally greater than 0.9). Let $w \in [0, 1]$ be an arbitrary ray.
    For each time-step $t$, we can find $u_{w, t}$ with $ \mathbb{P}(K_{w, t} > u_{w, t}  \ |\  Z_t = z_t) = 1 - q_1$, and analogously for $v_{w, t}$ and $q_2$.
    Because we have $1 - q_2 < 1 - q_1$, by monotonicity of CDF we have $v_{w, t} \geq u_{w, t}$ for each $t$. 
    Moreover, assuming strict monotonicity, we have $v_{w, t} > u_{w, t}$, and can use \cref{eq: cond lambda main} to get the estimate $\hat{\lambda}_q$, given in \cref{def: lambda qr}.
    \begin{align} 
      && e^{-(v_{w, t} - u_{w, t}) \lambda(w \ |\ z_t)} &= \mathbb{P}(K_{w, t} > v_{w, t} \ |\ K_{w, t} > u_{w, t}, Z_t = z_t) \\
      && &= \frac{\mathbb{P}(K_{w, t} > v_{w, t}  \ |\ , Z_t = z_t)}{\mathbb{P}(K_{w, t} > u_{w, t}  \ |\ , Z_t = z_t)}  \\
      && &= \frac{1 - q_2}{1 - q_1} \\
      \implies && \hat{\lambda}_q (w \ |\ z_t) &= -\frac{1}{v_{w, t} - u_{w, t}} \log \left(\frac{1 - q_2}{1 - q_1}\right) \label{def: lambda qr}
    \end{align}

    We now have a way to estimate $\lambda(w)$ through a sequence of thresholds $\{ u_{w, t}, v_{w, t}\}$ conditional on $z_t$ and corresponding quantiles. It remains to determine how to get the thresholds from the data.
    This is commonly done through quantile regression. 
    We use the standard assumption of linearity in the covariates to estimate $u_{w, t} = z_t^T \pi_u$, where $\pi_u \in \mathbb{R}^g$ is estimated by minimizing the expected loss $\argmin_{\pi \in \mathbb{R}^g} \sum_{i = 1}^n (K_{w, t} - z_t^T\pi)(q_i - \1 _{K_{w, t} < z_t^T\pi})$ (Koenker, 2017) and analogously for $v_{w, t}$.

    A major weakness of this method is that quantiles $q_1, q_2$ are modelled separately, so we cannot guarantee $v_{w, t} \geq u_{w, t}$. 
    To mitigate this issue, MBW propose to use GAMs from (Andre, 2023), modelling exceedences as a GPD, shown in \cref{eq: gpd biv}. 
    This is akin to the formulation in \cref{eq: gpd} with constant shape parameter and using link function $h(\tau(z_t)) = \log(\tau(z_t))$. 
    We take some base threshold $u_{w, t}^*$ correponding to quantile $q < q_1 < q_2$, estimated using the process described above. 
    We can then estimate $u_{w, t}, v_{w, t}$ from the fitted GPD model as its $\frac{q_1 - q}{1 - q}$ and $\frac{q_2 - q}{1 - q}$ quantiles respectively.
    MBW show through simulation that these GAM quantile estimates give better estimates of $\lambda(w)$ than standard quantile regression.
    \begin{equation} \label{eq: gpd biv}
      K_{w, t} - u_{w, t}^* \ |\ K_{w, t} > u_{w, t}^*, Z_t = z_t \sim \GP(\tau(z_t), \xi)
    \end{equation}

    Our estimates $\hat{\lambda}_q$ are dependent the quantiles $q_1, q_2$ chosen. In fact (EXPLANATION). Hence, MBW use a set of pairs of quantiles $\{ q_{1j}, q_{2j}\}_{j = 1}^m$ with each $q_{1j} < q_{2j} < 1, \; j=1,\dots,m$. 
    We can then take the estimator $\hat{\lambda}_{qave}$ as shown in \cref{def: lambda qr ave} by averaging over individual $\hat{\lambda}_q$ estimates. 
    MBW also show that this estimator outperforms estimators obtained by a single quantile pair.
    \begin{equation} \label{def: lambda qr ave}
      \hat{\lambda}_{qave}(w \ |\ z_t) = \frac{1}{m} \sum_{j = 1}^m \hat{\lambda}_q(w \ |\ z_t)
    \end{equation}

  \subsection{Bernstein-Bezier Polynomial Estimator}
    










  \subsection{Limitations}





\newpage
\section{test}

We made amazing contributions to the world of musical fractal pasta 
\citep{McDonald2017,Tibshirani2013}. We use Natbib, so be sure to use
\citep{Stein1981} for parenthetical references. Or you can say, according to
\citet{HastieTibshirani2009}, we should strive to balance truth and lies.


\bibliographystyle{rss}
\bibliography{qp.bib}

\end{document}
